\section{Soal - Simulasi Semifinal OSN SMP Pelatda DKI Jakarta 2025}
Diketahui $ABCD$ adalah segiempat talibusur dimana $AB$ dan $CD$ berpotongan di $E$. Titik $K$ dan $L$ terletak pada $CD$ dan $AB$ berturut-turut sehingga $\angle KAD = \angle KBC$ dan $\angle LDA = \angle LBC$. Buktikan bahwa $EK = EL$.

\definecolor{qqqqcc}{rgb}{0.,0.,0.8}
\definecolor{ffqqtt}{rgb}{1.,0.,0.2}
\definecolor{qqccqq}{rgb}{0.,0.8,0.}
\definecolor{wewdxt}{rgb}{0.43137254901960786,0.42745098039215684,0.45098039215686275}
\definecolor{rvwvcq}{rgb}{0.08235294117647059,0.396078431372549,0.7529411764705882}
\begin{center}
    \begin{tikzpicture}[line cap=round,line join=round,>=triangle 45,x=1.0cm,y=1.0cm,scale=0.5]
        \clip(-13,-12) rectangle (16,9);
        \draw [line width=1.pt,color=wewdxt] (-12.811872319363635,-4.793793271603575)-- (2.730245744242749,1.5178560751872407);
        \draw [line width=1.pt,color=wewdxt] (2.730245744242749,1.5178560751872407)-- (4.752108806286289,-7.909080451590773);
        \draw [line width=1.pt,color=wewdxt] (4.752108806286289,-7.909080451590773)-- (-12.811872319363635,-4.793793271603575);
        \draw [line width=1.pt,color=qqccqq] (0.3126582275437109,-3.9309513941472476) circle (5.96105965737421cm);
        \draw [line width=1.pt,color=wewdxt] (-5.219413690183677,-1.7104981665681662)-- (-5.2241310875286535,-6.139615315905112);
        \draw [line width=1.pt,color=wewdxt] (-5.219413690183677,-1.7104981665681662)-- (-1.2576596879951634,-6.843140183778444);
        \draw [line width=1.pt,color=wewdxt] (-1.2576596879951634,-6.843140183778444)-- (2.730245744242749,1.5178560751872407);
        \draw [line width=1.pt,color=wewdxt] (-5.2241310875286535,-6.139615315905112)-- (-1.9396346433208076,-0.3785809083134162);
        \draw [line width=1.pt,color=wewdxt] (-1.9396346433208076,-0.3785809083134162)-- (4.752108806286289,-7.909080451590773);
        \draw [line width=1.pt,color=wewdxt] (-5.214070383544707,3.306278953983742)-- (1.7492492096329906,6.091752417805242);
        \draw [line width=1.pt,color=ffqqtt] (-0.4279009035287008,-2.107361898627917) circle (4.807920091526327cm);
        \draw [line width=1.pt,color=wewdxt] (1.7492492096329906,6.091752417805242)-- (4.752108806286289,-7.909080451590773);
        \draw [line width=1.pt,color=wewdxt] (-5.214070383544707,3.306278953983742)-- (-5.2241310875286535,-6.139615315905112);
        \draw [line width=1.pt,color=wewdxt] (-5.214070383544707,3.306278953983742)-- (4.752108806286289,-7.909080451590773);
        \draw [line width=1.pt,color=wewdxt] (1.7492492096329906,6.091752417805242)-- (-5.2241310875286535,-6.139615315905112);
        \draw [line width=1.pt,color=wewdxt] (-1.9396346433208076,-0.3785809083134162)-- (-1.2576596879951634,-6.843140183778444);
        \draw [line width=1.pt,dash pattern=on 5pt off 5pt,color=qqqqcc] (0.753727638967831,-1.4442031845049497) circle (7.601426842221794cm);
        \draw [line width=1.pt,color=ffqqtt] (0.06322475363186655,-5.33725667991257) circle (5.347888477876674cm);
    
    \begin{scriptsize}
        \draw [fill=rvwvcq] (-12.811872319363635,-4.793793271603575) circle (2.5pt);
        \draw[color=rvwvcq] (-12.543464809954179,-4.0595015753729085) node {$E$};
        \draw [fill=rvwvcq] (2.730245744242749,1.5178560751872407) circle (2.5pt);
        \draw[color=rvwvcq] (2.993696203702673,2.2403920694680592) node {$B$};
        \draw [fill=rvwvcq] (4.752108806286289,-7.909080451590773) circle (2.5pt);
        \draw[color=rvwvcq] (5.0034782253697285,-7.190123570661979) node {$C$};
        \draw [fill=rvwvcq] (-5.219413690183677,-1.7104981665681662) circle (2.5pt);
        \draw[color=rvwvcq] (-4.968132574439893,-1.0061788886094944) node {$A$};
        \draw [fill=wewdxt] (-5.2241310875286535,-6.139615315905112) circle (2.0pt);
        \draw[color=wewdxt] (-4.968132574439893,-5.4895387830975455) node {$D$};
        \draw [fill=rvwvcq] (-1.2576596879951634,-6.843140183778444) circle (2.5pt);
        \draw[color=rvwvcq] (-0.98721818536861,-6.146582905565622) node {$K$};
        \draw [fill=wewdxt] (-1.9396346433208076,-0.3785809083134162) circle (2.0pt);
        \draw[color=wewdxt] (-1.6829119620995137,0.2692597020638302) node {$L$};
        \draw [fill=wewdxt] (-5.214070383544707,3.306278953983742) circle (2.0pt);
        \draw[color=wewdxt] (-4.9294829201770645,3.940976857032492) node {$N$};
        \draw [fill=wewdxt] (1.7492492096329906,6.091752417805242) circle (2.0pt);
        \draw[color=wewdxt] (2.027454847131973,6.72375196395611) node {$M$};
    \end{scriptsize}
    \end{tikzpicture}
\end{center}

\newpage
\section{Solusi}
Tanpa mengurangi keumuman, misalkan $A$ berada di antara $E$ dan $B$. Misalkan $M$ dan $N$ berturut-turut adalah perpotongan $BC$ dengan $DL$ dan $AD$ dengan $CL$.

\definecolor{qqqqcc}{rgb}{0.,0.,0.8}
\definecolor{ffqqtt}{rgb}{1.,0.,0.2}
\definecolor{qqccqq}{rgb}{0.,0.8,0.}
\definecolor{wewdxt}{rgb}{0.43137254901960786,0.42745098039215684,0.45098039215686275}
\definecolor{rvwvcq}{rgb}{0.08235294117647059,0.396078431372549,0.7529411764705882}
\begin{center}
    \begin{tikzpicture}[line cap=round,line join=round,>=triangle 45,x=1.0cm,y=1.0cm,scale=0.5]
        \clip(-13,-12) rectangle (16,9);
        \draw [line width=1.pt,color=wewdxt] (-12.811872319363635,-4.793793271603575)-- (2.730245744242749,1.5178560751872407);
        \draw [line width=1.pt,color=wewdxt] (2.730245744242749,1.5178560751872407)-- (4.752108806286289,-7.909080451590773);
        \draw [line width=1.pt,color=wewdxt] (4.752108806286289,-7.909080451590773)-- (-12.811872319363635,-4.793793271603575);
        \draw [line width=1.pt,color=qqccqq] (0.3126582275437109,-3.9309513941472476) circle (5.96105965737421cm);
        \draw [line width=1.pt,color=wewdxt] (-5.219413690183677,-1.7104981665681662)-- (-5.2241310875286535,-6.139615315905112);
        \draw [line width=1.pt,color=wewdxt] (-5.219413690183677,-1.7104981665681662)-- (-1.2576596879951634,-6.843140183778444);
        \draw [line width=1.pt,color=wewdxt] (-1.2576596879951634,-6.843140183778444)-- (2.730245744242749,1.5178560751872407);
        \draw [line width=1.pt,color=wewdxt] (-5.2241310875286535,-6.139615315905112)-- (-1.9396346433208076,-0.3785809083134162);
        \draw [line width=1.pt,color=wewdxt] (-1.9396346433208076,-0.3785809083134162)-- (4.752108806286289,-7.909080451590773);
        \draw [line width=1.pt,color=wewdxt] (-5.214070383544707,3.306278953983742)-- (1.7492492096329906,6.091752417805242);
        \draw [line width=1.pt,color=ffqqtt] (-0.4279009035287008,-2.107361898627917) circle (4.807920091526327cm);
        \draw [line width=1.pt,color=wewdxt] (1.7492492096329906,6.091752417805242)-- (4.752108806286289,-7.909080451590773);
        \draw [line width=1.pt,color=wewdxt] (-5.214070383544707,3.306278953983742)-- (-5.2241310875286535,-6.139615315905112);
        \draw [line width=1.pt,color=wewdxt] (-5.214070383544707,3.306278953983742)-- (4.752108806286289,-7.909080451590773);
        \draw [line width=1.pt,color=wewdxt] (1.7492492096329906,6.091752417805242)-- (-5.2241310875286535,-6.139615315905112);
        \draw [line width=1.pt,color=wewdxt] (-1.9396346433208076,-0.3785809083134162)-- (-1.2576596879951634,-6.843140183778444);
        \draw [line width=1.pt,dash pattern=on 5pt off 5pt,color=qqqqcc] (0.753727638967831,-1.4442031845049497) circle (7.601426842221794cm);
        \draw [line width=1.pt,color=ffqqtt] (0.06322475363186655,-5.33725667991257) circle (5.347888477876674cm);
    
    \begin{scriptsize}
        \draw [fill=rvwvcq] (-12.811872319363635,-4.793793271603575) circle (2.5pt);
        \draw[color=rvwvcq] (-12.543464809954179,-4.0595015753729085) node {$E$};
        \draw [fill=rvwvcq] (2.730245744242749,1.5178560751872407) circle (2.5pt);
        \draw[color=rvwvcq] (2.993696203702673,2.2403920694680592) node {$B$};
        \draw [fill=rvwvcq] (4.752108806286289,-7.909080451590773) circle (2.5pt);
        \draw[color=rvwvcq] (5.0034782253697285,-7.190123570661979) node {$C$};
        \draw [fill=rvwvcq] (-5.219413690183677,-1.7104981665681662) circle (2.5pt);
        \draw[color=rvwvcq] (-4.968132574439893,-1.0061788886094944) node {$A$};
        \draw [fill=wewdxt] (-5.2241310875286535,-6.139615315905112) circle (2.0pt);
        \draw[color=wewdxt] (-4.968132574439893,-5.4895387830975455) node {$D$};
        \draw [fill=rvwvcq] (-1.2576596879951634,-6.843140183778444) circle (2.5pt);
        \draw[color=rvwvcq] (-0.98721818536861,-6.146582905565622) node {$K$};
        \draw [fill=wewdxt] (-1.9396346433208076,-0.3785809083134162) circle (2.0pt);
        \draw[color=wewdxt] (-1.6829119620995137,0.2692597020638302) node {$L$};
        \draw [fill=wewdxt] (-5.214070383544707,3.306278953983742) circle (2.0pt);
        \draw[color=wewdxt] (-4.9294829201770645,3.940976857032492) node {$N$};
        \draw [fill=wewdxt] (1.7492492096329906,6.091752417805242) circle (2.0pt);
        \draw[color=wewdxt] (2.027454847131973,6.72375196395611) node {$M$};
    \end{scriptsize}
    \end{tikzpicture}
\end{center}

    \begin{lemmarev}Lingkaran luar $\triangle DLC$ menyinggung $AB$ di $L$ dan lingkaran luar $\triangle ABK$ menyinggung $CD$ di $K$

    \begin{proof} Perhatikan bahwa 
    $$\angle NDM = \angle ADL = \angle LCB = \angle NCM$$
    yang menunjukkan segiempat $MNDC$ siklis.

    Dari sini karena $ABCD$ dan $MNDC$ siklis, kita peroleh
    $$\angle ABC = 180\dg - \angle ADC = 180\dg - \angle NDC = \angle NMC$$
    yang menunjukkan bahwa $AB \parallel MN$.

    Oleh karena itu didapat
    $$\angle ALD = \angle NMD = \angle NCD = \angle LCD.$$
    Sehingga dari Alternate Segment Theorem, lingkaran luar $\triangle DLC$ menyinggung $AB$ di $L$. Analogi cara tersebut, didapat lingkaran luar $\triangle ABK$ menyinggung $CD$ di $K$.
    \end{proof}
    \end{lemmarev}

    Sekarang, dengan Power of a Point, dari Lemma dan lingkaran $(ABCD)$ kita punya 
    $$EL^2 = ED \cdot EC = EA \cdot AB = EK^2 \implies EL = EK.$$
    Soal terbukti.
